\uuid{IkUZ}
\exo7id{2911}
\titre{exo7 2911}
\auteur{quercia}
\organisation{exo7}
\datecreate{2010-03-08}
\isIndication{false}
\isCorrection{false}
\chapitre{Dénombrement}
\sousChapitre{Binôme de Newton et combinaison}
\cpge{mpsi-mp2i.calcul-algebrique-trigonometrie.a, pcsi.calcul-algebrique-trigonometrie.a, ptsi.calcul-algebrique-trigonometrie.a, bcpst1.methodes-calcul}
\module{Algèbre}
\niveau{L1}
\difficulte{}

\contenu{
\question{
Soit $u_n = \sum_{p=0}^n C_{n-p}^p$. Montrer que $u_0 = u_1 = 1$ et :
$\forall\ n \in \N,\ u_{n+2} = u_{n+1} + u_n$
({\it suite de Fibonacci\/}).
}
}
